\documentclass[10pt,a4paper]{amsart}
\usepackage[margin=19mm]{geometry}
\usepackage{amsmath,amssymb,mathtools}
\usepackage[scr=rsfs,cal=euler]{mathalfa}
\usepackage{xfrac}
\usepackage[unicode,hidelinks,bookmarks=false]{hyperref}
\newcommand{\ie}{i.e.\ }
\newcommand{\cf}{cf.\ }
\newcommand{\resp}{respectively\ }
\newcommand{\Subsection}{\subsection}
\providecommand{\nobreakdash}{\nobreak\hbox{-}}
\setlength{\emergencystretch}{2em}
\multlinegap0pt
\usepackage{bm}
\usepackage{bbm}
\usepackage{dsfont}
\usepackage{xspace}
\usepackage{cancel}
\usepackage{mathtools}
\usepackage{amsthm}
\makeatletter
\@ifundefined{theorem}{\newtheorem{theorem}{Theorem}}{}
\makeatother
\newcommand{\set}{\mathsf{set}}
\title[Ranking theories via encoded $\beta$-models]{Ranking theories via encoded $\beta$-models}
\author{Hanul Jeon, Patrick Lutz, Fedor Pakhomov, James Walsh}
\date{Source: arXiv:2503.20470v2 (CC BY 4.0)}
\begin{document}
\maketitle

\noindent\emph{Source and licence.} Ranking theories via encoded $\beta$-models. Version arXiv:2503.20470v2 (CC BY 4.0), distributed under the Creative Commons Attribution 4.0 International licence, as recorded in the arXiv licence metadata for this exact version and shown on the abstract page https://arxiv.org/abs/2503.20470v2. Full rights notice: licenses/arxiv-2503.20470-CC-BY-4.0.txt.\par\smallskip

\noindent\textit{Editorial scope.} Counted unit: the Theorem whose source label is \texttt{Theorem: Real encoding property of Z2} in the article (source0.tex lines 506--508, source label \texttt{Theorem: Real encoding property of Z2}), delivered together with the article's own complete original proof of it (source0.tex lines 509--531, one \texttt{proof} environment of 23 lines). No mathematics is written, repaired or re-derived: statement and proof are the article's own text line for line. Required context retained uncounted: Theorem: Groszek's theorem (lines 493--503). Every definition, display and label of those lines is retained unchanged. Omitted from this package and not redistributed: 1--492, 504--505, 532--884. Nothing inside a retained excerpt is omitted or altered: every retained body is delivered exactly as the article's lines are written, so no omission receipt of any kind accompanies this package. Citations: the retained excerpts carry no citation command at all, so no bibliography entry of the article is transcribed and no reference is redistributed. Reference adaptation: none was needed and none was made; every reference inside the delivered unit resolves inside it, so no unresolved reference or citation is produced. Printed slips kept literal: no printed word, symbol, hypothesis or number of the counted statement or of its proof was altered. Layout and class: the article's own class is \texttt{amsart} with packages including amssymb, graphicx, todonotes, hyperref, xcolor, breakurl, aliascnt; the wrapper is the lane's pinned \texttt{amsart} editorial wrapper at 10pt on A4, which restates the theorem-like environments the retained text needs and adds the article's own macro definitions that the retained text needs (\texttt{\textbackslash set}), with extra wrapper packages (bm, bbm, dsfont, xspace, cancel, mathtools). The delivered numbering is the unit's own. Assets: no retained span contains a figure environment, a graphics-inclusion command, a PDF-page inclusion, a table or a TikZ picture; no image file is copied into this package, the package declares no asset dependency and no image file is redistributed with it. Licence: version arXiv:2503.20470v2 of this article is distributed under the Creative Commons Attribution 4.0 International licence, as recorded in the arXiv licence metadata for this exact version and shown on the abstract page https://arxiv.org/abs/2503.20470v2. A full rights notice is delivered at licenses/arxiv-2503.20470-CC-BY-4.0.txt. Characters: every retained excerpt is pure ASCII, so the delivered engine, XeLaTeX with its default Latin Modern text font, reports no missing character.
\begin{theorem}[Groszek] \label{Theorem: Groszek's theorem} \pushQED{\qed}
    Let $X\subseteq \omega_1$, and suppose that there is maximum $a$ in the constructibility degree. Then there is a forcing $\mathbb{P}$ preserving $\omega_1$ such that for every $\mathbb{P}$-generic filter $G$ over $V$, the constructibility degrees above $a$ in $V[G]$ are precisely
    \begin{equation*}
        \{d_\gamma\mid \gamma<\omega_1 \} \cup \{a_\gamma,b_\gamma\mid \gamma\in X\}
    \end{equation*}
    satisfying the following:
    \begin{enumerate}
        \item For $\gamma<\beta$, $d_\gamma <_L d_\beta$.
        \item When $\gamma\in X$, $d_\gamma <_L a_\gamma,b_\gamma<_L d_{\gamma+1}$, and $a_\gamma$ and $b_\gamma$ are $\le_L$-incompatible. \qedhere 
    \end{enumerate}
\end{theorem}

\begin{theorem} \label{Theorem: Real encoding property of Z2}
    $\mathsf{Z}_2^\set$ has the real encoding property.
\end{theorem}

\begin{proof}
    For a given real $X$, consider the theory $T_X$ comprising the following axioms:
    \begin{enumerate}
        \item Axioms of $\mathsf{Z}_2^\set$.
        \item There is a real $Y$ such that the degree of constructibility above $Y$ forms a well-founded partial order of height $\omega+\omega$, and moreover
        \begin{enumerate}
            \item Each level of the constructibility degrees above $Y$ has at most two constructibility degrees.
            \item For every $n<\omega$, there is only one constructibility degree of level $n$ above $Y$.
            \item For each (meta-)natural number $n$, the $(\omega+h_n+n+1)$-th level of the construcbility degree over $Y$ has two incompatible degrees, where $h_n$ is the number of elements in $\{m<n\mid m\in X\}$.
            \item For (meta-)natural $m$ not of the form $h_n+n+1$, the $(\omega+m)$-th level of the construcbility degree over $Y$ has only one degree.
        \end{enumerate}
    \end{enumerate}
    $T_X$ requires that the degree of constructibility codes $\{\omega + n\mid n\in X\}$ above the top constructiblity degree. 
    We will construct a model of $T_X$ by forcing over $L[X]$ by applying \autoref{Theorem: Groszek's theorem}. $L[X]$ has $X$ as the top in the constructibility degree, so the reader may wonder why we do not simply add a constructibility degree pattern coding $X$ instead of padding the $\omega$-sequence of degrees first.
    However, the constructibility degrees over $L[X]$ may have a tail of constructibility degrees isomorphic to $-\omega$. In this case, the generic extension $L[X][G]$ may not find where the top real in the ground model is. To avoid losing the starting point, we add the $\omega$-sequence of reals and then code the real after the $\omega$-sequence, so that the generic extension can read off the real from the last limit point of the constructibility degree.

    Every transitive model of $T_X$ contains $X$ since $T_X$ can define $X$ from the constructibility degrees. Hence it remains to show that $T_X$ has a countable transitive model.
    Let $X$ be a real, and consider the inner model $L[X]$, which thinks there is a maximum real in the constructibility degree (namely $X$).
    Working in $L[X]$, let $G$ be a $\mathbb{P}$-generic filter over $L[X]$ for the forcing poset $\mathbb{P}$ given by \autoref{Theorem: Groszek's theorem} coding the set $\{\omega+n\mid n\in X\}$ into the constructibility degree. Then we can see that $M_0 = (H_{\omega_1})^{L[X][G]}$ is a model of $T_X$.
    
    However, we proved the existence of $M_0$ in $L[X][G]$ and not in $V$, and $M_0$ may not be countable. To solve both issues, let us work in $L[X][G]$ and consider a countable elementary submodel $M_1\prec M_0$. We can encode $M_1$ as a real, which we can treat as a well-founded model of $T_X$. Thus $L[X][G]$ satisfies the assertion ``There is a countable well-founded model of $T_X$,'' which is a $\Sigma^1_2$-assertion with a parameter $X$.
    By Shoenfield absoluteness, the assertion also holds in $V$, so $V$ thinks $T_X$ has a countable well-founded model. Collapsing a countable well-founded model of $T_X$ gives a desired model.
\end{proof}

\end{document}
